\section{The discovery: memoized inference in the KV cache}
\label{sec:mechanism}

\paragraph{Setup.} We study a minimal but agent-realistic decision: a context contains a policy rule and a
\emph{mutable field} whose value determines the correct action (e.g.\ ``cancel an order only if its status
is \texttt{pending}''; with \texttt{status=shipped} the correct action flips from \emph{cancel} to
\emph{deny}). After prefilling the full context we compare four cache states at the decision token:
\emph{stale} (old field, old downstream), \emph{field-only} (refresh the field's KV, leave the downstream
stale), \emph{full-downstream} (recompute everything after the field), and \emph{oracle} (a clean prefill
of the new value). We report \emph{decision recovery}: the fraction of the oracle's flip that a cache
state reproduces ($0$ = behaves like stale, $1$ = behaves like oracle). Four independent probes converge
on one account (\Cref{fig:mech}).

\paragraph{(1) Locality: the field's own KV barely matters.} Refreshing only the field's KV recovers
essentially none of the decision---field-only recovery is $-0.028$ on Llama-3.1-8B and near zero across
models---whereas recomputing the downstream recovers it fully ($1.0$) (\Cref{fig:mech}a). The field is
read \emph{indirectly}: its causal contribution to the decision is under $1\%$.

\paragraph{(2) Suffix concentration: the effect lives downstream and late.} Sweeping how many downstream
tokens we patch (ranked by causal effect) shows recovery accrues slowly and saturates only as we include
many tokens after the field, with the causal mass concentrated in mid/late layers
(\Cref{fig:mech}b). The information the decision needs is not in the field but distributed over the tokens
that followed it at prefill (\Cref{fig:mech}c).

\paragraph{(3) Linear decodability.} The field-conditioned conclusion is linearly decodable from those
downstream tokens' residual stream at prefill time---i.e.\ the model has already \emph{computed and
written down} the answer, not merely copied the field.

\paragraph{(4) Knockout and dose-response.} Ablating the high-effect downstream tokens flips the decision
back to stale, and the effect grows with the number/position of memoizing tokens, ruling out a diffuse
explanation: specific aggregator/delimiter positions carry the conclusion. This mirrors the delimiter-token
aggregation reported by \citet{lindsey2025biology} for forward planning; here the stored quantity is a
backward-looking, field-conditioned \emph{conclusion}.

\paragraph{What the note contains.} If the note were a verbatim copy of the field, restating the value
would suffice and emphatic wording would be inert. Instead, a wording ablation (\Cref{fig:mech}d) shows
that once the corrected value is present, the override phrasing is \emph{redundant} (bare value, tagged
update, and explicit override all reach $\approx 1.0$ P(safe)), while aggressive ``disregard your earlier
conclusion and re-evaluate'' phrasing actively \emph{hurts} ($0.81$). The note behaves like a committed
conclusion that a late, salient correction can overwrite---but that confrontational instructions can
destabilize. We name the phenomenon \textbf{attention-mediated memoized inference}, and the rest of the
paper shows it makes the cache both editable (\Cref{sec:editable}) and composable (\Cref{sec:composable}).

\begin{figure}[t]
\centering
\includegraphics[width=\linewidth]{fig2_mechanism}
\caption{\textbf{Four causal probes for memoized inference.} (a) Refreshing the field's own KV recovers
$\approx 0$ of the decision; recomputing the downstream recovers it fully. (b) Recovery is
suffix-concentrated, accruing only as many post-field tokens are patched. (c) The decision reads almost
entirely from downstream notes, not the field token. (d) Once the value is present, override wording is
redundant and ``re-evaluate'' phrasing hurts---the note is a committed conclusion.}
\label{fig:mech}
\end{figure}
